\documentclass[a4paper, 10pt, twocolumn]{article}
\include{amsymb}
\include{marvosym}

\usepackage[francais]{babel}
\usepackage[cyr]{aeguill}
\usepackage[applemac]{inputenc}
\usepackage{graphicx}
\usepackage{xspace}
\usepackage[a4paper]{geometry}
\usepackage{latexsym,amsmath,amssymb,textcomp}
\usepackage{moreverb}
\usepackage{listings}
\usepackage{multirow}
\usepackage{titling}
\usepackage{mathabx}

\usepackage{pdfpages}

\usepackage{hyperref}


\newlength{\indentationnota}
\newlength{\largeurlignenota}
\newlength{\paddingnota}
\newlength{\largeurnota}
\setlength{\paddingnota}{5pt}
\setlength{\largeurnota}{0.9cm}

\definecolor{pink}{rgb}{1,0.5,1}
\makeatletter

\newenvironment{pictonote}[1]{% on passe le nom du �chier en argument
\begin{list}{}{%
\setlength{\labelsep}{0pt}%
\setlength{\rightmargin}{15pt}%
\setlength{\paddingnota}{5pt}}
\item%
\setlength{\indentationnota}{%
\@totalleftmargin+\largeurnota+\paddingnota}%
\setlength{\largeurlignenota}{%
\linewidth-\largeurnota-\paddingnota}%
\parshape=3%
\indentationnota\largeurlignenota%
\indentationnota\largeurlignenota%
\@totalleftmargin\linewidth%
\raisebox{-\largeurnota+2.2ex}[0pt][0pt]{%
\makebox[0pt][r]{%
\includegraphics[width=\largeurnota]{#1}%
\hspace{\paddingnota}}}%
\ignorespaces}{%
\end{list}}

\makeatother

\newenvironment{attention}%
{\begin{pictonote}{/Users/benjamin/Documents/Education/LaTeX/danger}}%
{\end{pictonote}}

\newenvironment{unicorn}%
{\begin{pictonote}{/Users/benjamin/Documents/Education/LaTeX/angry_unicorn}}%
{\end{pictonote}}

\newcommand\unicornbox[1]{\colorbox{pink}{\color{white}{#1}}}

\newcommand\cf{\emph{c.f}}

\newcommand\signature{%
\begin{figure}[br]
	\begin{flushright}
	\begin{minipage}{8cm}
		\begin{center}
		\reflectbox{\includegraphics[width=8cm]{/Users/benjamin/Documents/Education/LaTeX/unicorn}}\\
		\theauthor
		\end{center}
	\end{minipage}
	\end{flushright}
\end{figure}}

\newcommand\danger{\raisebox{-0.4ex}{\LARGE $\triangle$ \normalsize} \hspace{-4.1ex}! \hspace{1ex}}
\newcommand\PL{Programmation Logique\xspace}
\newcommand\pl{\bsc{Prolog}\xspace}


\title{RDF - S�ance 4 : Les filtres de Laws}
\author{Benjamin \bsc{Van Ryseghem}}

\begin{document}
\maketitle

\section{Filtres de Laws}
\subsection{Question 1}
\subsubsection{Image noire}

\paragraph{Filtre 1}
On obtient que des $-16$.

\paragraph{Filtre 2}
On obtient que des $0$.

\paragraph{Filtre 3}
On obtient que des $0$.

\paragraph{Filtre 4}
On obtient que des $0$.

\paragraph{Filtre 5}
On obtient que des $0$.

\subsubsection{Image blanche}

\paragraph{Filtre 1}
On obtient que des $16$.

\paragraph{Filtre 2}
On obtient que des $0$.

\paragraph{Filtre 3}
On obtient que des $0$.

\paragraph{Filtre 4}
On obtient que des $0$.

\paragraph{Filtre 5}
On obtient que des $0$.

\subsection{Question 2}

\paragraph{Filtre 1}
On obtient que des $0$.

\paragraph{Filtre 2}
On obtient que des $0$.

\paragraph{Filtre 3}
On obtient que des $0$.

\paragraph{Filtre 4}
On obtient que des $0$.

\paragraph{Filtre 5}
On obtient que des $[-16,16,-16,16, \dots\ ]$.


\subsection{Question 3}

\paragraph{Filtre 1}
Il y a $0$ alternance.

\paragraph{Filtre 2}
Il y a $1$ alternance.

\paragraph{Filtre 3}
Il y a $2$ alternance..

\paragraph{Filtre 4}
Il y a $3$ alternance.

\paragraph{Filtre 5}
Il y a $4$ alternance.

\subsection{Question 4}

\paragraph{Filtre 1}
Il faut une ligne constante de $1$.

\paragraph{Filtre 2}
\dots,$-1$,$-1$,$-1$,$-1$,$1$,$1$,$1$,$1$, \dots

On obtient alors \dots, $6$, $2$, $-2$, $-6$, $-6$,$-2$,$2$,$6$, \dots

\paragraph{Filtre 3}
\dots, $-1$, $-1$, $1$, $1$, \dots

On obtient alors \dots, $-4$, $-4$, $4$, $4$, \dots

\paragraph{Filtre 4}
\dots , -1, -1,1,-1,1,1,-1,1 \dots

On obtient alors \dots, $6$, $-6$, $2$, $2$, $-6$, $6$, $-2$, $-2$ \dots

\paragraph{Filtre 5}
Il faut une alternance noir, blanc.

\onecolumn
\signature

\end{document}